% Options for packages loaded elsewhere
\PassOptionsToPackage{unicode}{hyperref}
\PassOptionsToPackage{hyphens}{url}
\documentclass[
  11pt,
]{article}
\usepackage{xcolor}
\usepackage[margin=1in]{geometry}
\usepackage{amsmath,amssymb}
\setcounter{secnumdepth}{-\maxdimen} % remove section numbering
\usepackage{iftex}
\ifPDFTeX
  \usepackage[T1]{fontenc}
  \usepackage[utf8]{inputenc}
  \usepackage{textcomp} % provide euro and other symbols
\else % if luatex or xetex
  \usepackage{unicode-math} % this also loads fontspec
  \defaultfontfeatures{Scale=MatchLowercase}
  \defaultfontfeatures[\rmfamily]{Ligatures=TeX,Scale=1}
\fi
\usepackage{lmodern}
\ifPDFTeX\else
  % xetex/luatex font selection
\fi
% Use upquote if available, for straight quotes in verbatim environments
\IfFileExists{upquote.sty}{\usepackage{upquote}}{}
\IfFileExists{microtype.sty}{% use microtype if available
  \usepackage[]{microtype}
  \UseMicrotypeSet[protrusion]{basicmath} % disable protrusion for tt fonts
}{}
\makeatletter
\@ifundefined{KOMAClassName}{% if non-KOMA class
  \IfFileExists{parskip.sty}{%
    \usepackage{parskip}
  }{% else
    \setlength{\parindent}{0pt}
    \setlength{\parskip}{6pt plus 2pt minus 1pt}}
}{% if KOMA class
  \KOMAoptions{parskip=half}}
\makeatother
\usepackage{graphicx}
\makeatletter
\newsavebox\pandoc@box
\newcommand*\pandocbounded[1]{% scales image to fit in text height/width
  \sbox\pandoc@box{#1}%
  \Gscale@div\@tempa{\textheight}{\dimexpr\ht\pandoc@box+\dp\pandoc@box\relax}%
  \Gscale@div\@tempb{\linewidth}{\wd\pandoc@box}%
  \ifdim\@tempb\p@<\@tempa\p@\let\@tempa\@tempb\fi% select the smaller of both
  \ifdim\@tempa\p@<\p@\scalebox{\@tempa}{\usebox\pandoc@box}%
  \else\usebox{\pandoc@box}%
  \fi%
}
% Set default figure placement to htbp
\def\fps@figure{htbp}
\makeatother
\setlength{\emergencystretch}{3em} % prevent overfull lines
\providecommand{\tightlist}{%
  \setlength{\itemsep}{0pt}\setlength{\parskip}{0pt}}
\usepackage{bookmark}
\IfFileExists{xurl.sty}{\usepackage{xurl}}{} % add URL line breaks if available
\urlstyle{same}
\hypersetup{
  pdftitle={ACT-2008 -- Rapport d'étape 2},
  pdfauthor={Équipe 17 : Joël Ducasse, Alexandre Quirion, Ahlin Roland Césaire Tchintchin et Marc-Antoine Dion},
  hidelinks,
  pdfcreator={LaTeX via pandoc}}

\title{ACT-2008 -- Rapport d'étape 2}
\author{Équipe 17 : Joël Ducasse, Alexandre Quirion, Ahlin Roland
Césaire Tchintchin et Marc-Antoine Dion}
\date{Le 26 Mars 2026}

\begin{document}
\maketitle

\section*{1. Mandat}

Le mandat de cette seconde étape était d'avancer notre application
Shiny. Plus précisément, nous avions vu un deuxième modèle de prévision
en classe. Il s'agit du modèle de Bühlmann. Le mandat principal était
d'implanter ce deuxième modèle dans notre application Shiny. De plus,
nous devions commencer à réfléchir sur le troisième et dernier modèle de
crédibilité à implanter.

\section*{2. Réalisations}

Lors de cette deuxième étape, nous avons principalement implanté le
modèle de Bühlmann dans notre application Shiny. Tout d'abord, pour
effectuer avec succès l'implantation de ce modèle, nous avons commencé
par estimer les différents paramètres du modèle de Bühlmann.
C'est-à-dire que nous avons estimé les paramètres \(s^2\), \(a\) et
\(m\). Pour les paramètres \(s^2\) et \(a\) ̂, nous avons utilisé les
formules de la méthode Bühlmann et les données du travail pratique pour
estimer ces paramètres. Plus précisément, nous avons utilisé les données
complètes dans le fichier « resultats-complets.csv ». Pour le paramètre
m, nous avons, aussi, utilisé ce fichier de données complètes. Par
contre, nous avons aussi normalisé les résultats à l'aide des données
des pars (« normales.csv »). Plus précisément, nous avons soustrait les
pars aux coups réalisés par les golfeurs pour tenir compte de la
difficulté variable des trous dans le terrain de golf. Par la suite, on
a calculé la valeur de \(K\) et de \(Z\) à l'aide des valeurs
précédemment calculées. De plus, nous avons programmé dans notre
application le calcul de la moyenne individuelle S. Enfin, la dernière
étape était de programmer la crédibilité Bühlmann.

Nous avons effectué quelques tests et pour le moment tout semble
cohérent. En effet, un golfeur qui a joué plusieurs trou aura un facteur
de crédibilité Z plus élevé qu'un golfeur ayant réalisé peu de coup.
Selon nous, ceci est logique et concorde avec le modèle de Bühlmann
puisque lorsqu'une personne a plus de données, on a une plus grande
confiance en son expérience individuelle. Donc, les tests étaient
cohérents, mais une petite révision de ce modèle avant la remise serait
bénéfique pour être certain que tout concorde.

Voici les formules que nous avons utilisé pour les différents paramètres
et pour la crédibilité Bühlmann lors de l'implantation de ce modèle :

\[
m = E[\mu(\Theta_i)];
\qquad
s^2 = E[\sigma^2(\Theta_i)];
\qquad
a = \mathrm{Var}(\mu(\Theta_i))
\]

\[
Z = \frac{n}{n + K};
\qquad
K = \frac{s^2}{a}
\] \[
\Pi^{B}_{i,n+1} = Z\,\bar{S}_i + (1 - Z)\,m;
\qquad
\bar{S}_i = \frac{1}{n}\sum_{t=1}^{n} S_{it}
\]

\section*{3. Apprentissages}

Lors de cette seconde étape, nous avons appris à bien coder le modèle de
Bühlmann dans une application Shiny. En effet, on a réussi à coder
chacun des paramètres de ce modèle et d'approximer les futurs coups des
golfeurs à travers ce modèles.

De plus, nous avons appris un nouveau modèle (Bühlmann-Straud) en
classe. Ceci nous sera fort utile pour la troisième étape de ce projet.
En effet, il s'agit du troisième et dernier modèle à implanter dans
notre application.

\section*{4. Comparaison du temps de réalisation de cette seconde étape}

Lors du café du savoir 2, nous avions estimé que cette seconde étape
allait être d'une durée d'environ 6 à 8 heures. Notre temps réel
travaillé sur l'implantation de la méthode Bühlman au sein de notre
application était autour de 8 heures. Notre estimation était donc assez
réelle puisque nous sommes dans l'intervalle prédit. Nous avons donc eu
assez de temps pour réaliser cette étapes et nous somme donc prêt à
réaliser l'implantation du dernier modèle dans notre application. La
plus longue étapes de cette étape a été d'implanter les formules des
différents paramètres du modèle de Bühlman.

\end{document}
